\section{Introduction} \label{sec:intro}
The earliest hints of the existence of Dark Matter (DM) came from the dynamics of stars and galaxies \citep{zwicky:1933,rubin:1970}. 
These observations provided the first evidence that galaxies are surrounded by DM halos extending more than an order of magnitude past the visible matter \citep{1974ApJ...193L...1O,1974Natur.250..309E}. In the Milky Way, the DM halo dominates the mass of the Galaxy, and is expected to be $\sim$ 1-2 orders of magnitudes larger in mass than the baryonic (stars and gas) component. However, many of its detailed properties have thus far not been well understood. 
Insight into both the total mass and density profile of the DM halo is extremely valuable for galactic dynamics and studies of DM, informing our understanding of baryonic feedback \citep{2016MNRAS.456.3542T,2020MNRAS.497.2393L}, the particle nature of DM \citep{Spergel:1999mh,Tulin:2017ara,Lisanti:2018qam,Nadler:2019zrb,DES:2020fxi}, and indirect detection of DM \citep{Cirelli:2010xx,Slatyer:2017sev,Rinchiuso:2020skh}.  


Many methods have been used to measure the MW mass, and to a lesser extent the density profile of DM: the circular velocity of the MW \citep{nesti:2013, deSalas:2019, eilers:2019,Ou:2023,2023arXiv230900048J,2023ApJ...942...12W,2023ApJ...946...73Z}, the population and motion of satellites \citep{barber:2014, cautun:2014, patel:2018, callingham:2019, sohn:2018, watkins:2019, fritz:2020}, stellar streams \citep{gibbons:2014, kupper:2015, dierickx:2017, Vasiliev:2021}, and the escape velocity of the MW \citep{Piffl2014, williams:2017, banik:2018, Monari2018, deason:2019, Koppelman2021, Necib2022, Prudil:2022}. The escape velocity, in particular, gives a measure of the gravitational potential at a given location, and can be used to constrain the DM potential when combined with models of the baryonic mass. It is typically extracted by modeling the tail of the stellar speed distribution. Due to a lack of data, earlier studies focused on the local escape velocity in the solar neighborhood \citep{Leonard1990,smith:2007}. 
In this paper, we produce precise measurements of the escape velocity over a range of Galactocentric radii, in order to obtain an escape velocity profile and thus inform the DM potential over a wider range. We revisit prior power law models for the tail of the speed distribution, and also introduce a new, more robust functional form. 


Our analysis is made possible with the significant improvements in the third data release from \Gaia~(\Gaia~DR3). The \Gaia space mission \citep{gaia:2016, gaiaDR1, gaiaDR2, GaiaeDR3} has revolutionized the field of MW astronomy, from studies of streams \citep{malhan:2018, price-whelan:2018, helmi:2020, chandra:2023}, dust \citep{green:2019, lallement:2019, leike:2020}, disk resonances \citep{kawata:2018, trick:2021}, to the discovery of previously unknown merging events \citep{2018Natur.563...85H,2018MNRAS.478..611B,2020Necib,Naidu2020}, and more. The latest data release of the \Gaia mission includes positions, parallax, proper motions, and radial velocities of an astounding 33~million stars \citep{gaiaDR3, gaiaRVs}. The six-dimensional features of this large dataset allow us to finally probe detailed dynamics of the MW at a few~kpc from the solar position, enabling more precise and reliable determinations of the escape velocity profile. \Gaia~DR3, in particular, has a significant advantage over DR2 in that it includes many more sources at large distances from the Sun with line-of-sight velocities. 

Even with excellent data, estimating the escape velocity crucially relies on modeling the tail of the stellar speed distribution.  \cite{Leonard1990} introduced a power-law model, with a distribution given by
\begin{equation} \label{eq:powerlaw}
    g(v \mid \vesc, k) \propto (\vesc - v)^{k}, \quad v\in[\vmin,\vesc]
\end{equation}
where $\vesc$ is the escape velocity, $k$ is the slope of the power law, and $\vmin$ is a lower speed cutoff which reflects the fact that we only model the \textit{tail} of the speed distribution for stars at a given Galactocentric radius. The overall normalization is set such that $\int_{\vmin}^{\vesc} dv \, g(v) = 1$. This functional form has been the basis of most of the escape velocity measurements in the MW for the last few decades. However, escape velocity estimates based upon this modeling \citep{smith:2007, Piffl2014, Monari2018, deason:2019, Koppelman2021, Necib2022} face various challenges such as small sample sizes, degeneracy of parameters, systematic uncertainties in model selection, and the choice of strong prior distributions on the fitting parameters. 

The functional form  in Eq.~\ref{eq:powerlaw} exhibits a strong degeneracy between the two parameters $\vesc$ and $k$, leading to large uncertainties on the escape velocity estimates; this degeneracy can be exacerbated by small sample sizes or if the underlying distribution deviates from a single power law. To reduce such uncertainties on $\vesc$, previous work in the literature adopted an array of either theoretically or numerically (based on cosmological simulations) motivated priors on the slope $k$. Such choices affect the final measurement of the escape velocity, and ideally one would require a model that best fits the data without the need to adopt strong priors. 

\begin{figure*}[t!]
  \centering
  \includegraphics[width=\textwidth]{media/dataset_and_fits.pdf}
  \caption{High-speed and high-quality sample of stars in \Gaia~DR3 with 6D kinematics, and the derived escape velocity profile at $68\%$ (boxes) and $95\%$ (whiskers) confidence intervals. Medians marked with a horizontal line. Escape velocities shown here are obtained by binning the data into $1\,\rm{kpc}$ bins and fitting the speed distributions above $\vmin = 310\,\rm{km\,s^{-1}}$ with the stretched exponential power law model introduced in this paper. Contours correspond to the logarithm of the number of stars at that speed and Galactocentric radius, labelled N.}
  \label{fig:data}
\end{figure*}

Motivated by the discovery of the \Gaia-Sausage-Enceladus (GSE) \citep{2018Natur.563...85H,2018MNRAS.478..611B} merger that has contributed a large fraction of the ex-situ stars in the MW \citep[see e.g.][]{Necib:2019a,naidu:2020}, and arguing for the need to address the issues above, \cite{Necib2022methods} adopted a two-component power law, generalizing Eq.~\ref{eq:powerlaw} but with two different slopes $k_1$ and $k_2$, 
\begin{align} \label{eq:powerlaw2}
    g(v \mid \vesc, k_1, k_2) \propto &(1-f_S) g(v \mid \vesc, k_1)   \\ 
    & + f_S g(v \mid \vesc, k_2), \quad v\in[\vmin,\vesc] \nonumber
\end{align}
where $f_S$ is an additional parameter that controls the relative fraction of the power laws. Based on this new functional form, \cite{Necib2022} computed the escape velocity of the MW at the position of the Sun using the early \Gaia data release 3 (eDR3) \citep{GaiaeDR3}. 

The two-component model in \cite{Necib2022methods} provided a more general functional form, allowing for a better fit to data containing substructure. Indeed, \cite{Necib2022} showed using the Akaike information criterion (AIC; \cite{AIC}), that the two-component model is preferred at lower $\vmin$, but that at higher $\vmin$ it was not distinguishable from a single component. This affirms the presence of a deviation from the single power law away from $\vesc$. The model further addressed the degeneracy and the issue of strong priors on the slope $k$: with the modified distribution, it was possible to obtain a good fit to the data and robust measurement of $\vesc$ with wide priors on the slopes.

In this work, we aim at further investigating models for the tail of the speed distribution, adapting to the improved quality and increased statistics of the data in \Gaia DR3.  We reconsider the single power law and  two-component power law models, and obtain the first two-component power law escape velocity profile using this extended dataset. With the improved dataset, we find that there remains a number of issues in using these power-law models: they can predict a high $\vesc$ relative to the fastest stars in the dataset, results can be sensitive to the choice of $\vmin$, and some parameters correlated with the escape velocity can remain unconverged even with wide priors. Motivated by these challenges, we introduce a new functional form that allows for a steeper rise in the speed distribution away from $\vesc$, while keeping a power law cutoff near $\vesc$. We find that this new form most closely tracks the fastest stars and gives the most stable results. With each of these models, we obtain independent escape velocity measurements in 1~kpc-wide bins from 4~kpc to 11~kpc in Galactocentric radius. We use the resulting escape velocity profile to fit the DM density profile, and combine this with a circular velocity measurement from \cite{eilers:2019} to obtain an updated measurement of the MW mass.

The remainder of this paper is organized as follows: In Section \ref{sec:data}, we present the selection criteria for the high-quality sample of \Gaia~DR3 data used to obtain our escape velocity profile and Milky Way halo constraints. In Section \ref{sec:methods} we describe the general modeling approach for the tail of the stellar speed distribution, and in Section \ref{sec:modeling} we present the different models for this speed distribution. These models include the previously studied one and two-component power laws, as well as the new model we introduce in this work, the ``stretched exponential power law." In Section \ref{sec:results}, we show the results of applying these three models to the \Gaia~DR3 data, presenting the escape velocity as a function of the radial distance and validating the modeling. Finally in Section \ref{sec:DM}, we use the escape velocity data to constrain the DM density profile and total mass of the Milky Way.